% Source: https://artofproblemsolving.com/community/c3h1248599_symmetry_of_a_cube
% Stable identity: aops-topic-1248599-q01
A plane of symmetry of a cube passes through its centre and maps the cube to itself. There are two kinds.

\begin{asy}
usepackage("amsmath");
size(3.5cm);
defaultpen(fontsize(10pt));
import three; currentprojection=orthographic(4,2.5,2);
draw(unitbox,gray);
draw(surface((0.5,0,0)--(0.5,1,0)--(0.5,1,1)--(0.5,0,1)--cycle),lightblue+opacity(0.6));
draw(surface((0,0,0)--(1,1,0)--(1,1,1)--(0,0,1)--cycle),pink+opacity(0.6));
label("\mbox{}",(min(currentpicture,user=true)+max(currentpicture,user=true))/2,black); // keeps the text after the figure black when LaTeX typesets the labels inline
\end{asy}

\begin{itemize}
\item Planes parallel to a pair of faces, through the midpoints of the four edges perpendicular to them (blue): one for each pair of opposite faces, $3$ planes.
\item Planes through two opposite edges, containing a diagonal of two opposite faces (pink): there are $12$ edges forming $6$ pairs of opposite edges, giving $6$ planes.
\end{itemize}
No other plane works (a symmetry plane must map vertices to vertices and so must be one of these). In total $3+6=\boxed{9}$ planes.
